% Build from the repository root (auxiliary files stay under build/):
% mkdir -p build/camera-tree-writeup
% latexmk -pdf -interaction=nonstopmode -halt-on-error \
%   -outdir=build/camera-tree-writeup docs/technical/camera_tree_intrinsics.tex
% cp build/camera-tree-writeup/camera_tree_intrinsics.pdf docs/technical/
\documentclass[11pt,letterpaper]{article}
\usepackage[T1]{fontenc}
\usepackage{lmodern}
\usepackage[margin=0.9in]{geometry}
\usepackage{amsmath,amssymb,bm,mathtools}
\usepackage{booktabs,tabularx}
\usepackage{microtype}
\usepackage{xcolor}
\usepackage{enumitem}
\usepackage{fancyhdr}
\usepackage{listings}
\usepackage[colorlinks=true,linkcolor=blue!45!black,urlcolor=blue!45!black]{hyperref}
\hypersetup{pdftitle={Retaining Intrinsics in Camera-Tree Schur Complements},pdfauthor={GTSFM technical note}}
\pagestyle{fancy}
\fancyhf{}
\fancyhead[L]{\small GTSFM technical note}
\fancyhead[R]{\small Camera-tree merging}
\fancyfoot[C]{\thepage}
\setlength{\headheight}{14pt}
\setlength{\parindent}{0pt}
\setlength{\parskip}{5pt plus 1pt}
\setlist{nosep,leftmargin=*}
\lstset{basicstyle=\ttfamily\small,columns=fullflexible,keepspaces=true,
  frame=single,rulecolor=\color{black!20},breaklines=true,aboveskip=8pt,belowskip=8pt}
\newcommand{\R}{\mathbb{R}}
\newcommand{\SE}{\mathrm{SE}(3)}
\newcommand{\Sim}{\mathrm{Sim}(3)}
\newcommand{\Local}{\operatorname{Local}}
\newcommand{\rank}{\operatorname{rank}}
\newcommand{\diag}{\operatorname{diag}}
\newcommand{\argmin}{\operatorname*{arg\,min}}
\newcommand{\code}[1]{\texttt{\detokenize{#1}}}
\begin{document}
\thispagestyle{empty}
\begin{center}
{\LARGE\bfseries Retaining Intrinsics in\\[4pt] Camera-Tree Schur Complements\par}
\vspace{9pt}
{\large A technical description of the GTSFM implementation}\\[5pt]
{\small November 12, 2026 \quad $\cdot$ \quad PR \href{https://github.com/borglab/gtsfm/pull/1188}{\#1188}}
\end{center}
\vspace{3pt}
\begin{abstract}
\noindent Independently reconstructed image clusters can be merged by combining
frozen, landmark-free camera likelihoods. To refine calibration during merging,
the reduced state must contain both poses and intrinsics: only landmark increments
are eliminated. This note derives the retained quadratic, its gauge treatment,
and its conversion into similarity-aware factors for a single global solve.
It also specifies the public data contract, fixed-intrinsics compatibility mode,
and the limits of the resulting approximation.
\end{abstract}

\section{State, observations, and the intended approximation}
A leaf reconstruction supplies camera poses, intrinsics, image measurements, and
3D landmarks in an arbitrary local frame. Distinct leaves can contain the same
image, identified by filename, but use unrelated local keys and coordinate frames.
The merging problem introduces one global pose and calibration for each image
and one similarity transformation for each leaf.

\begin{center}
\begin{tabularx}{\linewidth}{@{}lX@{}}
\toprule
Symbol & Meaning \\
\midrule
$\ell$, $i$, $j$ & Leaf, globally identified image, and leaf-local landmark indices. \\
$\bar T_{\ell i}\in\SE$ & Frozen camera-to-leaf pose estimate. \\
$\bar\kappa_{\ell i}\in\R^3$ & Frozen \code{Cal3Bundler} parameters $(f,k_1,k_2)$. \\
$\bar p_{\ell j}\in\R^3$ & Frozen landmark position in the leaf frame. \\
$T_i,\kappa_i$ & Global camera pose and per-image calibration to optimize. \\
$S_\ell\in\Sim$ & Similarity mapping global coordinates into leaf $\ell$. \\
\bottomrule
\end{tabularx}
\end{center}

The observation model uses one focal length and two radial-distortion coefficients.
For $q=T^{-1}p$, let $x=q_x/q_z$, $y=q_y/q_z$, and $r^2=x^2+y^2$. Then
\begin{equation}
 \pi(T,\kappa,p)=
 \begin{bmatrix} f(1+k_1r^2+k_2r^4)x+u_0\\
                 f(1+k_1r^2+k_2r^4)y+v_0\end{bmatrix}.
\end{equation}
The principal point $(u_0,v_0)$ is fixed. Occurrences of a shared image must agree
on it when intrinsics are variables. Different images have independent
calibrations; this API does not tie parameters by physical camera identity.

The producer linearizes each retained observation once at the supplied leaf
estimates, with isotropic pixel noise $\sigma^2 I_2$ (default $\sigma=1$).
It does not run bundle adjustment first. Tracks with fewer than two distinct
retained cameras are excluded. The consumer subsequently optimizes this frozen
camera approximation, without reading landmarks or re-evaluating reprojections.
This reconstruction-based producer rebuilds a measurement model; it does not
recover the original BA graph, robust weighting, priors, or gauge from saved geometry.

\newpage
\section{Eliminating landmarks while retaining intrinsics}
For a leaf, collect pose increments $\delta x\in\R^{6n}$, calibration increments
$\delta\kappa\in\R^{3n}$, and landmark increments $\delta p\in\R^{3m}$.
Whitened linearization gives the local least-squares problem
\begin{equation}\label{eq:linear}
 \min_{\delta x,\delta\kappa,\delta p}
 \frac12\left\|A_x\delta x+A_\kappa\delta\kappa+A_p\delta p-b\right\|^2,
 \qquad b=-r(\bar T,\bar\kappa,\bar p).
\end{equation}
The residual offset $b$ is retained even when the supplied reconstruction is not
at a stationary point. Set $\delta c=[\delta x^\top\;\delta\kappa^\top]^\top$
and $A_c=[A_x\;A_\kappa]$. The normal equations have block structure
\begin{equation}
 \begin{bmatrix}H_{cc}&H_{cp}\\H_{pc}&H_{pp}\end{bmatrix}
 \begin{bmatrix}\delta c\\\delta p\end{bmatrix}
 =\begin{bmatrix}g_c\\g_p\end{bmatrix},\qquad
 H=A^\top A,\quad g=A^\top b.
\end{equation}
If $H_{pp}$ is nonsingular, profiling out the points yields
\begin{align}
 H_{\mathrm{red}}&=H_{cc}-H_{cp}H_{pp}^{-1}H_{pc},\label{eq:schur}\\
 g_{\mathrm{red}}&=g_c-H_{cp}H_{pp}^{-1}g_p.\label{eq:linear-term}
\end{align}
The pose--intrinsics block is therefore
\begin{equation}
 (H_{\mathrm{red}})_{x\kappa}
 =H_{x\kappa}-H_{xp}H_{pp}^{-1}H_{p\kappa}.
\end{equation}
Preserving this block allows calibration changes to influence the pose solution
and preserves correlations between images induced by shared landmarks.

\subsection{Square-root reduction used in the code}
The implementation uses QR elimination rather than explicitly forming
Equation~\eqref{eq:schur}. Conceptually, an orthogonal transformation gives
\begin{equation}
 Q^\top\begin{bmatrix}A_p&A_c&b\end{bmatrix}
 =\begin{bmatrix}R_p&B_p&d_p\\0&B&d\end{bmatrix}.
\end{equation}
When the landmark columns are eliminable, solving the first block for $\delta p$
leaves the reduced cost
\begin{equation}\label{eq:reduced}
 E_{\mathrm{red}}(\delta c)=\tfrac12\|B\delta c-d\|^2,
 \qquad B^\top B=H_{\mathrm{red}},\quad B^\top d=g_{\mathrm{red}}.
\end{equation}
Any additive constant discarded by the factor representation does not affect
optimization. GTSAM's \code{GaussianFactorGraph.marginal(retained_keys, EliminateQR)}
retains the pose and calibration keys and eliminates only landmark keys.
The returned producer summary is ungauged.

\subsection{Fixing, retaining, and marginalizing calibration differ}
Fixing intrinsics imposes $\delta\kappa=0$ before reduction, producing a
$6n$-dimensional camera system. Retaining intrinsics produces the $9n$-dimensional
system above. Marginalizing intrinsics would instead minimize over
$\delta\kappa$ and discard their variables; that is not the new behavior.
Likewise, an information matrix alone cannot reconstruct the reduced likelihood:
$g_{\mathrm{red}}$ is essential when $b\ne0$. For example, $x^2$ and $(x-3)^2$
have the same quadratic coefficient but different optima.

\newpage
\section{Gauge constraints and leaf Bayes trees}
An observable monocular leaf has seven similarity gauge freedoms, even when
its intrinsics are retained. For $n$ cameras, the expected ungauged rank is
$9n-7$ in variable-intrinsics mode and $6n-7$ in fixed-intrinsics mode.
Additional self-calibration or geometric ambiguities are rejected; the code does
not add implicit calibration priors to make such a system invertible.

The consumer fixes a local coordinate chart before factoring an ungauged leaf.
It adds six exact rows fixing the first pose at its frozen estimate and one exact
row fixing the baseline length to the farthest camera, linearized at those
estimates. A nonzero baseline is required. These constraints act only on poses,
not on calibration variables. GTSAM's \code{NonlinearEqualityPose3} and constrained
\code{RangeFactorPose3} supply the rows.

The augmented graph is factored by multifrontal QR with GTSAM's default ordering.
Its Gaussian Bayes-tree conditionals jointly represent the constrained camera
quadratic. A conditional can contain any mixture of pose and calibration blocks;
block widths are six and three, respectively. Conditionals preserve their linear
terms and noise models, including exact constraints.

External producers may supply already gauged graphs with \code{gauge_fixed=True}.
These must be full rank, and the consumer does not add a second coordinate gauge.
Mixed gauged and ungauged leaves are supported. Rank checks normalize Jacobian
columns and tolerate numerical perturbations of gauge null modes; rank alone
cannot certify that an external producer's constraints are purely coordinate
choices. That remains part of the producer's contract.

\section{Lifting the leaf quadratic to global variables}
A similarity variable lets each leaf's constrained local chart move within the
global reconstruction. Let $\Phi(S,T)$ denote GTSAM's
\code{Similarity3.transformFrom(Pose3)}. With GTSAM's stored parameters
$S=(R_s,t_s,s)$, its action is
\begin{equation}
 p\longmapsto s(R_sp+t_s),\qquad
 \Phi(S,(R,t))=(R_sR,\;s(R_st+t_s)).
\end{equation}
The translation parameter is inside the scale in this convention. The code uses
the library operation and its derivatives, rather than implementing an alternate
similarity convention.

For each image in leaf $\ell$, define frozen-chart increments
\begin{align}
 \xi_{\ell i}(T_i,S_\ell)
 &=\Local_{\bar T_{\ell i}}\!\left(\Phi(S_\ell,T_i)\right)\in\R^6,\\
 \eta_{\ell i}(\kappa_i)
 &=\Local_{\bar\kappa_{\ell i}}(\kappa_i)
 =\kappa_i-\bar\kappa_{\ell i}\in\R^3.
\end{align}
For conditional $q$, concatenate the appropriate increments in its actual key
order into $z_{\ell q}$. The lifted residual is
\begin{equation}\label{eq:lift}
 e_{\ell q}=A_{\ell q}z_{\ell q}-b_{\ell q}.
\end{equation}
Each shared image has one global calibration but may have different frozen
calibration estimates in different leaves. Similarities affect poses only:
neither focal length nor distortion coefficients change under a change of world
frame or world scale. The residual is passed unwhitened to a \code{CustomFactor}
with the conditional's original noise model, so weighting is applied exactly once.

\newpage
\section{Derivatives, root optimization, and the final reduction}
The lifted factors use analytic derivatives in GTSAM's local retraction charts.
For a pose block, let $A_i^x$ be its six columns in a conditional and define
\begin{equation}
 J_i^{\mathrm{loc}}=D_2\Local_{\bar T_{\ell i}}(\Phi(S_\ell,T_i)),\quad
 J_i^S=D_S\Phi(S_\ell,T_i),\quad J_i^T=D_T\Phi(S_\ell,T_i).
\end{equation}
The chain rule gives
\begin{align}
 \frac{\partial e_{\ell q}}{\partial \delta T_i}
   &=A_i^xJ_i^{\mathrm{loc}}J_i^T,\\
 \frac{\partial e_{\ell q}}{\partial \delta S_\ell}
   &=\sum_{i\text{ in pose blocks}}A_i^xJ_i^{\mathrm{loc}}J_i^S,\\
 \frac{\partial e_{\ell q}}{\partial \delta\kappa_i}
   &=A_i^\kappa.
\end{align}
Calibration contributes no similarity derivative. In particular, a
calibration-only conditional has an identically zero similarity Jacobian.
The implementation computes column offsets from the variable types; it never
assumes that every conditional block has six columns. Finite-difference tests
check the derivatives at nonidentity poses and nonunit similarity scales.

\subsection{Initialization and one global solve}
Camera identity is resolved by filename, and sorted filenames determine the
output indices. Pose-based similarity alignment initializes a connected graph of
overlapping leaves. Each usable overlap link needs at least two shared cameras
with distinct centers. Global calibrations are initialized from the first
occurrence of each image in leaf traversal order.

The input scene tree assembles the leaf evidence bottom-up, without performing
an optimization at internal nodes. The root problem is
\begin{equation}\label{eq:root}
 \min_{\{T_i,\kappa_i,S_\ell\}}\ \sum_{\ell,q}
   \tfrac12\|e_{\ell q}(T,\kappa,S_\ell)\|_{\Sigma_{\ell q}}^2,
 \qquad S_0=I,
\end{equation}
Here $\|e\|_\Sigma^2=e^\top\Sigma^{-1}e$ for stochastic rows; zero-noise rows
are exact constraints. Fixing the first leaf's similarity selects the global frame and
scale. Levenberg--Marquardt with multifrontal QR solves this problem once at the
root. The local pose and calibration estimates used inside the factors remain
frozen throughout optimization. The public leaf-to-global transform is
$S_\ell^{-1}$.

\subsection{Final camera information retains calibration}
At the converged solution, the assembled graph is linearized again with respect
to its global camera and similarity variables. This relinearizes the lifted
factors, not the original image measurements. The final marginalization removes
only the similarity increments:
\begin{equation}
 E_{\mathrm{final}}(\delta c)
 =\min_{\delta s} E_{\mathrm{root,lin}}(\delta c,\delta s),\qquad
 \delta c=(\delta x,\delta\kappa).
\end{equation}
The resulting Gaussian Bayes tree contains global \code{X(i)} pose keys and
\code{K(i)} calibration keys. It contains no landmarks or similarities and
includes the chosen coordinate gauge. The returned scene has optimized poses
and calibrations, image metadata, and zero tracks. Its quadratic describes
local information in the frozen approximation, not the uncertainty of a newly
optimized full bundle-adjustment problem.

\newpage
\section{Public contract, compatibility, and validation}
The reconstruction wrapper and the external-factor route agree when supplied
the same reduced likelihoods. The reconstruction producer uses \code{GeneralSFMFactor2Cal3Bundler} to
linearize pose, point, and calibration variables. In fixed-intrinsics mode it
uses \code{GeneralSFMFactorPoseCal3Bundler} instead.

\begin{lstlisting}[language=Python]
# Default: eliminate points; retain and optimize pose + intrinsics.
result = optimize_camera_tree(scene_tree, CameraTreeOptions())

# Compatibility: keep intrinsics fixed before linearization.
fixed = optimize_camera_tree(
    scene_tree, CameraTreeOptions(optimize_intrinsics=False))
\end{lstlisting}

\code{CameraLeafFactors} carries the following contract:
\begin{itemize}
\item \code{camera_graph} is the complete reduced Gaussian likelihood, including
      its linear term and weighting; Jacobian and Hessian factors are accepted.
\item \code{estimated_camera_poses} and \code{camera_metadata} use the same local
      pose keys. Metadata contains the filename, frozen calibration, and image shape.
\item A nonempty \code{calibration_keys} maps every pose key to one distinct
      calibration key, disjoint from the pose keys. An empty mapping means fixed
      intrinsics and preserves existing constructor usage.
\item Graph keys must exactly equal the declared pose and calibration keys.
      Key symbols do not determine types. Arbitrary keys may be reused across leaves.
\item \code{gauge_fixed} declares whether the seven coordinate freedoms have
      already been removed. The consumer validates rank accordingly.
\end{itemize}

The consumer infers the variable layout from the summaries; producer options
\code{optimize_intrinsics} and \code{measurement_sigma} do not reinterpret or
reweight external factors. Entire pose-only trees remain supported. Trees mixing
pose-only and variable-intrinsics summaries are rejected. Copies are rekeyed to
global \code{X(i)} and \code{K(i)} without modifying caller-owned input data.
Leaf summaries expose frozen \code{estimated_calibrations}; output cameras expose
the optimized values.

\subsection{The original-graph marginal API is currently pose-only}
The current PR also includes \code{compute_camera_marginal_factors(graph, values)}
from PR \href{https://github.com/borglab/gtsfm/pull/1187}{\#1187}. This route
linearizes the actual BA factors, including their gauge, priors, and weights.
It eliminates non-camera variables first, then uses METIS ordering and QR to
produce compact pose conditionals. Its retained keys must contain \code{Pose3}
values: calibration variables, if present, are marginalized out. It does not
implement the joint retention described in Sections 2--5, and
\code{optimize_intrinsics} does not configure this API. Preserving original BA
factors and retaining intrinsics together requires an external joint summary
satisfying the contract above; extending this pose-only helper is separate work.

\newpage
\section{Validation, assumptions, and limitations}
At implementation commit \code{cae83707}, the optimizer and related tree test
suite reported \textbf{69 passed and one existing dataset-related skip} with
Python 3.12 and GTSAM 4.3.0. Coverage includes explicit point minimization versus
the retained quadratic, pose--intrinsics cross terms, conditional factorization,
weighted analytic derivatives, calibration reconciliation, and both intrinsics
modes on the overlapping 12-camera Lund fixture. It also checks frame invariance,
external-input equivalence, reused keys, Hessian factors, immutable inputs,
and failures for invalid or unobservable inputs. Local NumPy/Accelerate matrix
multiplication warnings persisted; explicit finite-result and equivalence checks
passed. These are correctness checks, not accuracy or performance benchmarks.

The main numerical checks target distinct parts of the derivation:
\begin{enumerate}
\item \textbf{Landmark profiling.} Compare reduced objective differences against
explicit least-squares minimization of point increments at nonzero residuals.
Check that pose--intrinsics cross terms survive and that factorization preserves
the constrained quadratic.
\item \textbf{Lifted factors.} Compare pose, calibration, and similarity
Jacobians with finite differences, including calibration-only conditionals.
Check original weighting against the Gaussian factor's cost.
\item \textbf{Merged outputs.} Reconcile independent calibration likelihoods
with known optima, verify exported intrinsics and retained keys, and compare
reconstruction-wrapper results with external-summary results.
\item \textbf{Coordinate and input invariants.} Change leaf frames independently,
reuse arbitrary local keys, verify caller inputs remain unchanged, and exercise
both existing gauge modes and both intrinsics modes.
\end{enumerate}

The reduction is exact for the chosen linearized least-squares problem,
up to discarded constants, but not for the original nonlinear BA objective.
The root solver re-evaluates pose-chart and similarity transformations while
keeping the measurement-derived quadratic frozen. It neither updates points nor
relinearizes image observations. Large changes from the frozen estimates can
therefore reduce approximation accuracy, even if the root objective decreases.

Leaf likelihoods must not count the same evidence multiple times. The Lund
fixture assigns each source track to exactly one leaf while allowing cameras to
overlap; the general API does not deduplicate observations. Shared calibration
across different images is also outside its contract. Choosing image identity
by filename assumes filenames are unique within a leaf and consistent across
leaves.

Observability and coordinate conventions are explicit preconditions, not
regularization choices. The usable overlap graph must be connected, its links
must provide nondegenerate baselines, and each leaf must have the declared rank.
Fixed-intrinsics observability does not guarantee self-calibration observability.
Principal points must agree across occurrences of a shared image. The final
Bayes tree includes the selected gauge; covariance interpretation must respect
that choice and the frozen measurement approximation.

The implementation remains separate from production scene merging. Serialization,
distributed execution, and further refinement against the original measurements
are outside this note's scope. The reconstruction convenience wrapper warns that
it cannot preserve the original BA gauge and noise models. The original-graph
helper avoids that loss but currently has the pose-only retention boundary
described in Section 6.1.

\vspace{3pt}
{\small\textbf{Implementation references.}
\href{https://github.com/borglab/gtsfm/blob/cae83707f52e92353c38a5db988b40def88fd2f5/gtsfm/camera_tree_optimizer.py}{Optimizer source};
\href{https://github.com/borglab/gtsfm/blob/cae83707f52e92353c38a5db988b40def88fd2f5/tests/test_camera_tree_optimizer.py}{numerical tests};
\href{https://github.com/borglab/gtsfm/pull/1184}{original PR \#1184};
\href{https://github.com/borglab/gtsfm/pull/1187}{original-graph API PR \#1187};
\href{https://github.com/borglab/gtsfm/pull/1188}{intrinsics PR \#1188}.
Joint-retention derivations refer to commit \code{cae83707}; the original-graph
API boundary was checked at \code{20649a17}. The accompanying repository Markdown
guide gives usage examples.}
\end{document}
